% !PRINTED 128
\weberpage{128}
%<<CONT
\[
-i\omega = r\,e^{i\varphi}, \qquad \alpha-\frac{\beta}{\omega} = \varrho\,e^{i\psi}
\]
\[
-i(-\beta+\alpha\omega) = r\varrho\,e^{i(\varphi+\psi)},
\]
then, since the real parts of $-i\omega$, $-i(-\beta+\alpha\omega)$ are positive,
\[
-\frac{\pi}{2} < \varphi < \frac{\pi}{2}, \quad -\pi < \psi < \pi, \quad
-\frac{\pi}{2} < \varphi+\psi < \frac{\pi}{2},
\]
\[
\sqrt{-i\omega} = \sqrt{r}\,e^{\frac{i\varphi}{2}}, \qquad
\sqrt{\alpha-\frac{\beta}{\omega}} = \sqrt{\varrho}\,e^{\frac{i\psi}{2}},
\]
and the real part of
\[
\sqrt{-i\omega}\ \sqrt{\alpha-\frac{\beta}{\omega}} = \sqrt{-i(-\beta+\alpha\omega)}
= \sqrt{r\varrho}\,e^{\frac{i(\varphi+\psi)}{2}}
\]
is positive. Accordingly there results from (15), 1:
\[
\sqrt{-i\omega}\, E\!\left(\begin{matrix}\alpha, & \beta \\ \gamma, & \delta\end{matrix}, \frac{-1}{\omega}\right)
= \left(\frac{\beta}{\alpha}\right) i^{\frac{\alpha-1}{2}}
e^{\frac{\pi i}{12}\left[\alpha(\gamma-\beta)-(\alpha^2-1)\beta\delta\right]}
\sqrt{-i(-\beta+\alpha\omega)}.
\]

The same formula, however, is obtained, since $\left(\dfrac{-\beta}{\alpha}\right)=i^{(\alpha-1)}\left(\dfrac{\beta}{\alpha}\right)$,
from (15), 2. for
\[
E\!\left(\begin{matrix}-\beta, & \alpha \\ -\delta, & \gamma\end{matrix}, \omega\right),
\]
and with that (14) is proved.

There still remains formula (13). It is enough to take account of the
upper signs alone, that is, to prove the formula
\[
E\!\left(\begin{matrix}\alpha+\beta, & \beta \\ \gamma+\delta, & \delta\end{matrix}, \omega\right)
= e^{\frac{\pi i}{12}} E\!\left(\begin{matrix}\alpha, & \beta \\ \gamma, & \delta\end{matrix}, \omega+1\right)
\]
since the other case reduces to this one by the interchange of $\alpha, \gamma, \omega$
with $\alpha+\beta, \gamma+\delta, \omega+1$. If in the first place
$\beta$ is odd (and positive), then (13) follows from (15), 2. on the
basis of the congruence
\[
(\beta^2-1)(1-\beta\gamma-\alpha\delta-\beta\delta) \equiv -\beta(\beta^2-1)(2\gamma+\delta) \equiv 0\ (\mathrm{mod}\ 24).
\]

If $\beta$ is even, then $\alpha$ is odd. If we take $\alpha$ positive, then
$\alpha+\beta$ can be positive or negative. If in the first case the
upper signs hold, in the second the lower, then (15), 1 gives us
\weq{16}{
e^{\frac{i\pi}{12}} E\!\left(\begin{matrix}\alpha, & \beta \\ \gamma, & \delta\end{matrix}, \omega+1\right)
= \left(\frac{\beta}{\alpha}\right) i^{-\frac{1-\alpha}{2}}
e^{\frac{\pi i}{12}\left[1+\alpha(\gamma-\beta)-(\alpha^2-1)\beta\delta\right]}
\sqrt{\alpha+\beta+\beta\omega},
}
%CONT>>
